\documentclass[11pt]{article}
\usepackage{../../latex/stat4ml_style}

\title{Chapter 2: Estimation and Losses \\ \large Part III --- Statistical Foundations of Deep Learning \& Reinforcement Learning}
\author{Omid Safarzadeh}
\date{}

\begin{document}
\maketitle

\begin{abstract}
\noindent
Every loss a network minimises is a likelihood in disguise, and every
regulariser is a prior. Squared error is the negative log-likelihood of Gaussian
noise, cross-entropy is the negative log-likelihood of a categorical label, and
once you see losses this way the choice of loss stops being arbitrary --- it is a
claim about the noise. The same move explains regularisation: adding a penalty is
multiplying the likelihood by a prior and maximising the posterior instead, with
the $L_2$ penalty a Gaussian prior and the $L_1$ penalty a Laplace prior whose
sharp peak at zero is exactly why lasso gives sparse weights. Threading both is
the \emph{score function}, the gradient of the log-likelihood, whose expectation
is zero at the truth --- the quantity estimation, and later policy gradients,
are built on.
\end{abstract}

\tableofcontents

% ======================================================================
\section{Losses are negative log-likelihoods}
% ======================================================================

A probabilistic model says the data were drawn from $p(y\mid x;\theta)$. The
\keyterm{maximum likelihood} estimate maximises $\prod_i p(y_i\mid x_i;\theta)$,
or equivalently minimises the \keyterm{negative log-likelihood}
$-\sum_i \log p(y_i\mid x_i;\theta)$. Different noise models give different
losses --- and they are the losses deep learning already uses.

\begin{theorem}[Squared error is Gaussian MLE]
If $y = f_\theta(x) + \varepsilon$ with $\varepsilon\sim N(0,\sigma^2)$, then the
negative log-likelihood is
\[
  -\sum_i \log p(y_i\mid x_i)
   = \frac{n}{2}\log(2\pi\sigma^2) + \frac{1}{2\sigma^2}\sum_i\big(y_i-f_\theta(x_i)\big)^2.
\]
Only the last term depends on $\theta$, so maximising the likelihood is exactly
minimising the sum of squared errors.
\end{theorem}
\begin{proof}
Substitute the Gaussian density $p(y_i\mid x_i)=
(2\pi\sigma^2)^{-1/2}\exp(-(y_i-f_\theta(x_i))^2/2\sigma^2)$, take logs, and sum.
The $\theta$-free terms are constant in the optimisation, leaving the SSE scaled
by $1/2\sigma^2$.
\end{proof}

\noindent The left panel of the figure makes the point concretely: the sum of
squared errors and the Gaussian negative log-likelihood, plotted against a
candidate slope, are the \emph{same} curve up to an affine transform, with the
same minimiser. The identical argument with a categorical likelihood gives
\keyterm{cross-entropy}: $-\sum_i \log p_{\text{model}}(y_i\mid x_i)$ with
$p_{\text{model}}$ the softmax output is the multinomial negative log-likelihood
(Part~II, Chapter~8). The loss is never arbitrary --- it names the noise.

% ======================================================================
\section{Regularisation is a prior}
% ======================================================================

Maximum likelihood overfits when the model is flexible. The Bayesian repair is a
\keyterm{prior} $p(\theta)$: instead of maximising the likelihood, maximise the
posterior $p(\theta\mid \text{data}) \propto p(\text{data}\mid\theta)\,p(\theta)$
--- \keyterm{maximum a posteriori} (MAP) estimation. Taking logs turns the product
into a sum, so MAP minimises
\[
  \underbrace{-\sum_i \log p(y_i\mid x_i;\theta)}_{\text{loss}}
  \;\; \underbrace{-\,\log p(\theta)}_{\text{penalty}}.
\]
The penalty \emph{is} the negative log-prior.

\begin{theorem}[Ridge is a Gaussian prior]
A Gaussian prior $\theta\sim N(0,\tau^2 I)$ contributes
$-\log p(\theta) = \tfrac{1}{2\tau^2}\|\theta\|^2 + \text{const}$. Combined with
the Gaussian likelihood of Theorem~1.1, MAP minimises
$\tfrac{1}{2\sigma^2}\sum_i(y_i-f_\theta(x_i))^2 + \tfrac{1}{2\tau^2}\|\theta\|^2$,
which is ridge regression with $\lambda = \sigma^2/\tau^2$.
\end{theorem}
\begin{proof}
The log of $p(\theta)\propto \exp(-\|\theta\|^2/2\tau^2)$ is
$-\|\theta\|^2/2\tau^2$ up to a constant. Add it to the negative log-likelihood
and multiply through by $2\sigma^2$: the objective is
$\sum_i(y_i-f_\theta(x_i))^2 + (\sigma^2/\tau^2)\|\theta\|^2$, the ridge objective
with $\lambda=\sigma^2/\tau^2$.
\end{proof}

Swap the Gaussian prior for a \keyterm{Laplace prior}
$p(\theta)\propto \exp(-\|\theta\|_1/b)$ and the penalty becomes
$\|\theta\|_1$ --- the lasso. The middle panel shows why the two behave
differently: the Laplace density has a sharp peak at zero, so its MAP estimate is
pulled onto zero for any weight the data do not strongly support, giving
\keyterm{sparsity}; the smooth Gaussian only shrinks. The right panel is the
geometric version --- the $L_1$ ``ball'' is a diamond, and the squared-error
contours first touch it at a corner, where a coordinate is exactly zero.

% ======================================================================
\section{The score function}
% ======================================================================

Both sections rest on one object: the \keyterm{score}, the gradient of the
log-likelihood with respect to the parameter,
\[
  s(\theta) = \nabla_\theta \log p(y\mid x;\theta).
\]
The MLE is where the total score vanishes, $\sum_i s(\theta)=0$, and the two
facts that make the score central are that, at the true parameter, its
expectation is zero and its variance is the \keyterm{Fisher information}.

\begin{theorem}[The score has mean zero]
Under regularity conditions, $\mathbb{E}_{y\sim p(\cdot;\theta)}[s(\theta)] = 0$.
\end{theorem}
\begin{proof}
$\mathbb{E}[s(\theta)] = \int \big(\nabla_\theta \log p\big)\,p\,dy
 = \int \frac{\nabla_\theta p}{p}\,p\,dy = \nabla_\theta\!\int p\,dy
 = \nabla_\theta 1 = 0$, using $\nabla_\theta\log p = \nabla_\theta p / p$ and
swapping gradient and integral.
\end{proof}

For the Gaussian mean, $s(\mu)=(y-\mu)/\sigma^2$, and $\sum_i s(\mu)=0$ gives
$\hat\mu = \bar y$ --- the sample mean is the score equation's solution. This is
not a curiosity: the identity $\nabla_\theta\log p = (\nabla_\theta p)/p$ is the
\keyterm{log-derivative trick}, and it returns in Chapter~7 as the engine of the
policy-gradient estimator, where the ``likelihood'' is a policy and the score is
differentiated through sampled actions.

\begin{summarybox}
A loss is a negative log-likelihood: squared error is the NLL of Gaussian noise,
cross-entropy the NLL of a categorical label, so the choice of loss is a claim
about the noise. A regulariser is a negative log-prior: MAP minimises
loss~$+$~penalty, with the $L_2$ penalty a Gaussian prior
($\lambda=\sigma^2/\tau^2$) and the $L_1$ penalty a Laplace prior whose peak at
zero yields sparsity. Both rest on the score $s(\theta)=\nabla_\theta\log p$,
whose expectation is zero at the truth (Fisher information is its variance), and
whose log-derivative identity $\nabla\log p=\nabla p/p$ reappears as the
policy-gradient estimator in Chapter~7.
\end{summarybox}

\end{document}
